\begin{textsample}{POS Dim 1 – 2024 – Score 20.00 – 2024/t003227.txt}  \label{ex:f1_pos_048}
Chris Anderson : Demis, so good to have you here.
Demis Hassabis : It ' s \textbf{fantastic} to be here, thanks, Chris.
Now, you told Time Magazine,"I want to understand \textbf{the} big questions, \textbf{the} really big ones that you normally go into philosophy or physics if you ' re interested in them.
I thought building AI would be \textbf{the} fastest route to answer some of those questions."Why did you think that?
DH : Well, I guess when I was a kid, my favorite subject was physics, and I was interested in all \textbf{the} big questions, fundamental nature of reality, what is consciousness, you know, all \textbf{the} big ones.
And usually you go into physics, if you ' re interested in that.
But I read a lot of \textbf{the} great physicists, some of my all-time scientific heroes like Feynman and so on.
And I realized, in \textbf{the} last, sort of 20, 30 years, we haven ' t made much progress in understanding some of these fundamental laws.
So I thought, why not build \textbf{the} ultimate tool to help us, which is \textbf{artificial} intelligence.
And at \textbf{the} same time, we could also maybe better understand ourselves and \textbf{the} brain better, by doing that too.
So not only was it incredible tool, it was also useful for some of \textbf{the} big questions itself.
CA : Super interesting.
So obviously AI can do so many things, but I think for this conversation, I ' d love to focus in on this \textbf{theme} of what it might do to \textbf{unlock} \textbf{the} really big questions, \textbf{the} giant scientific breakthroughs, because it ' s been such a \textbf{theme} driving you and your company.
DH : So I mean, one of \textbf{the} big things AI can do, and I ' ve always thought about, is we ' re getting, you know, even back 20, 30 years ago, \textbf{the} beginning of \textbf{the} internet era and \textbf{computer} era, \textbf{the} amount of data that was being produced and also scientific data, just too much for \textbf{the} human mind to comprehend in many cases.
And I think one of \textbf{the} uses of AI is to find patterns and insights in huge amounts of data and then \textbf{surface} that to \textbf{the} human scientists to make sense of and make new hypotheses and conjectures.
So it seems to me very compatible with \textbf{the} scientific method.
CA : Right.
But game play has played a huge role in your own journey in figuring this thing out.
Who is this young lad on \textbf{the} \textbf{left} there?
Who is that?
DH : So that was me, I think I must have been about around nine years old.
I ' m captaining \textbf{the} England Under 11 team, and we ' re playing in a Four Nations tournament, that ' s why we ' re all in red.
I think we ' re playing France, Scotland and Wales, I think it was.
CA : That is so weird, because that happened to me too.
In my dreams.
And it wasn ' t just chess, you loved all kinds of games.
DH : I loved all kinds of games, yeah.
CA : And when you launched DeepMind, pretty quickly, you started having it tackle game play.
Why?
DH : Well, look, I mean, games actually got me into AI in \textbf{the} first place because while we were doing things like, we used to go on training camps with \textbf{the} England team and so on.
And actually back then, I guess it was in \textbf{the} mid ' 80s, we would use \textbf{the} very early chess \textbf{computers}, if you remember them, to train against, as well as playing against each other.
And they were big lumps of plastic, you know, physical boards that you used to, some of you remember, used to actually press \textbf{the} squares down and there were LED lights, came on.
And I remember actually, not just thinking about \textbf{the} chess, I was actually just fascinated by \textbf{the} fact that this lump of plastic, someone had programmed it to be smart and actually play chess to a really high standard.
And I was just amazed by that.
And that got me thinking about thinking.
And how does \textbf{the} brain come up with these thought processes, these ideas, and then maybe how we could mimic that with \textbf{computers}.
So yeah, it ' s been a whole \textbf{theme} for my whole life, really.
CA : But you raised all this money to launch DeepMind, and pretty soon you were using it to do, for example, this.
I mean, this is an odd use of it.
What was going on here?
DH : Well, we started off with games at \textbf{the} beginning of DeepMind.
This was back in 2010, so this is from about 10 years ago, it was our first big breakthrough.
Because we started off with classic Atari games from \textbf{the} 1970s, \textbf{the} simplest kind of \textbf{computer} games there are out there.
And one of \textbf{the} reasons we used games is they ' re very convenient to test out your ideas and your algorithms.
They ' re really fast to test.
And also, as your systems get more powerful, you can choose harder and harder games.
And this was actually \textbf{the} first time ever that our machine surprised us, \textbf{the} first of many times, which, it figured out in this game called Breakout, that you could send \textbf{the} ball round \textbf{the} back of \textbf{the} wall, and actually, it would be much safer way to knock out all \textbf{the} tiles of \textbf{the} wall.
It ' s a classic Atari game there.
And that was our first real aha moment.
CA : So this thing was not programmed to have any strategy.
It was just told, try and figure out a way of winning.
You just move \textbf{the} bat at \textbf{the} bottom and see if you can find a way of winning.
DH : It was a real revolution at \textbf{the} time.
So this was in 2012, 2013 where we coined these terms"deep reinforcement \textbf{learning}."And \textbf{the} key thing about them is that those systems were learning directly from \textbf{the} pixels, \textbf{the} raw pixels on \textbf{the} screen, but they weren ' t being told anything else.
So they were being told, maximize \textbf{the} score, here are \textbf{the} pixels on \textbf{the} screen, 30, 000 pixels. \textbf{The} system has to make sense on its own from first principles what ' s going on, what it ' s controlling, how to get points.
And that ' s \textbf{the} other nice thing about using games to begin with.
They have clear objectives, to win, to get scores.
So you can kind of measure very easily that your systems are improving.
CA : But there was a direct line from that to this moment a few years later, where country of South Korea and many other parts of Asia and in fact \textbf{the} world went crazy over -- over what?
DH : Yeah, so this was \textbf{the} pinnacle of -- this is in 2016 -- \textbf{the} pinnacle of our games-playing work, where, so we ' d done Atari, we ' d done some more complicated games.
And then we reached \textbf{the} pinnacle, which was \textbf{the} game of Go, which is what they play in Asia instead of chess, but it ' s actually more complex than chess.
And \textbf{the} actual brute force algorithms that were used to kind of crack chess were not possible with Go because it ' s a much more pattern-based game, much more intuitive game.
So even though Deep Blue beat Garry Kasparov in \textbf{the} ' 90s, it took another 20 years for our program, AlphaGo, to beat \textbf{the} world champion at Go.
And we always thought, myself and \textbf{the} people working on this project for many years, if you could build a system that could beat \textbf{the} world champion at Go, it would have had to have done something very interesting.
And in this case, what we did with AlphaGo, is it basically learned for itself, by playing millions and millions of games against itself, ideas about Go, \textbf{the} right strategies.
And in fact invented its own new strategies that \textbf{the} Go world had never seen before, even though we ' ve played Go for more than, you know, 2, 000 years, it ' s \textbf{the} oldest board game in existence.
So, you know, it was pretty astounding.
Not only did it win \textbf{the} match, it also came up with brand new strategies.
CA : And you continued this with a new strategy of not even really teaching it anything about Go, but just setting up systems that just from first principles would play so that they could teach themselves from scratch, Go or chess.
Talk about AlphaZero and \textbf{the} amazing thing that happened in chess then.
DH : So following this, we started with AlphaGo by giving it all of \textbf{the} human games that are being played on \textbf{the} internet.
So it started that as a basic starting point for its knowledge.
And then we wanted to see what would happen if we started from scratch, from literally random play.
So this is what AlphaZero was.
That ' s why it ' s \textbf{the} zero in \textbf{the} name, because it started with zero prior knowledge And \textbf{the} reason we did that is because then we would build a system that was more general.
So AlphaGo could only play Go, but AlphaZero could play any two-player game, and it did it by playing initially randomly and then slowly, incrementally improving.
Well, not very slowly, actually, within \textbf{the} course of 24 hours, going from random to better than world-champion level.
CA : And so this is so amazing to me.
So I ' m more familiar with chess than with Go.
And for decades, thousands and thousands of AI experts worked on building incredible chess \textbf{computers}.
Eventually, they got better than humans.
You had a moment a few years ago, where in nine hours, AlphaZero taught itself to play chess better than any of those systems ever did.
Talk about that.
DH : It was a pretty incredible moment, actually.
So we set it going on chess.
And as you said, there ' s this rich history of chess and AI where there are these expert systems that have been programmed with these chess ideas, chess algorithms.
And you have this amazing, you know, I remember this day very clearly, where you sort of sit down with \textbf{the} system starting off random, you know, in \textbf{the} morning, you go for a cup of coffee, you come back.
I can still just about beat it by lunchtime, maybe just about.
And then you let it go for another four hours.
And by dinner, it ' s \textbf{the} greatest chess-playing entity that ' s ever existed.
And, you know, it ' s quite amazing, like, looking at that live on something that you know well, you know, like chess, and you ' re expert in and actually just seeing that in front of your eyes.
And then you extrapolate to what it could then do in science or something else, which of course, games were only a means to an end.
They were never \textbf{the} end in themselves.
They were just \textbf{the} training ground for our ideas and to make quick progress in a matter of, you know, less than five years actually went from Atari to Go.
CA : I mean, this is why people are in awe of AI and also kind of terrified by it.
I mean, it ' s not just incremental improvement. \textbf{The} fact that in a few hours you can achieve what millions of humans over centuries have not been able to achieve.
That gives you pause for thought.
DH : It does, I mean, it ' s a hugely powerful technology.
It ' s going to be incredibly transformative.
And we have to be very thoughtful about how we use that capability.
CA : So talk about this use of it because this is again, this is another extension of \textbf{the} work you ' ve done, where now you ' re turning it to something incredibly useful for \textbf{the} world.
What are all \textbf{the} letters on \textbf{the} \textbf{left}, and what ' s on \textbf{the} right?
DH : This was always my aim with AI from a kid, which is to use it to accelerate scientific discovery.
And actually, ever since doing my undergrad at Cambridge, I had this problem in mind one day for AI, it ' s called \textbf{the} protein-folding problem.
And it ' s kind of like a 50-year grand challenge in \textbf{biology}.
And it ' s very simple to explain.
Proteins are essential to life.
They ' re \textbf{the} building blocks of life.
Everything in your body depends on proteins.
A protein is sort of described by its amino acid sequence, which you can think of as roughly \textbf{the} genetic sequence describing \textbf{the} protein, so that are \textbf{the} letters.
CA : And each of those letters represents in itself a complex molecule?
DH : That ' s right, each of those letters is an amino acid.
And you can think of them as a kind of string of beads there at \textbf{the} bottom, left, right?
But in nature, in your body or in an animal, this string, a sequence, turns into this beautiful shape on \textbf{the} right.
That ' s \textbf{the} protein.
Those letters describe that shape.
And that ' s what it looks like in nature.
And \textbf{the} important thing about that 3D structure is \textbf{the} 3D structure of \textbf{the} protein goes a long way to telling you what its function is in \textbf{the} body, what it does.
And so \textbf{the} protein-folding problem is : Can you directly predict \textbf{the} 3D structure just from \textbf{the} amino acid sequence?
So literally if you give \textbf{the} machine, \textbf{the} AI system, \textbf{the} letters on \textbf{the} \textbf{left}, can it produce \textbf{the} 3D structure on \textbf{the} right?
And that ' s what AlphaFold does, our program does.
CA : It ' s not \textbf{calculating} it from \textbf{the} letters, it ' s looking at patterns of other folded proteins that are known about and somehow learning from those patterns that this may be \textbf{the} way to do this?
DH : So when we started this project, actually straight after AlphaGo, I thought we were ready.
Once we ' d cracked Go, I felt we were finally ready after, you know, almost 20 years of working on this stuff to actually tackle some scientific problems, including protein folding.
And what we start with is painstakingly, over \textbf{the} last 40-plus years, experimental biologists have pieced together around 150, 000 protein structures using very complicated, you know, X-ray crystallography \textbf{techniques} and other complicated experimental \textbf{techniques}.
And \textbf{the} rule of thumb is that it takes one PhD student their whole PhD, so four or five years, to uncover one structure.
But there are 200 million proteins known to nature.
So you could just, you know, take forever to do that.
And so we managed to actually fold, using AlphaFold, in one year, all those 200 million proteins known to science.
So that ' s a billion years of PhD time saved.
CA : So it ' s amazing to me just how reliably it works.
I mean, this shows, you know, here ' s \textbf{the} model and you do \textbf{the} experiment.
And sure enough, \textbf{the} protein turns out \textbf{the} same way.
Times 200 million.
DH : And \textbf{the} more deeply you go into proteins, you just start appreciating how exquisite they are.
I mean, look at how beautiful these proteins are.
And each of these things do a special function in nature.
And they ' re almost like works of art.
And it ' s still astounds me today that AlphaFold can predict, \textbf{the} green is \textbf{the} ground truth, and \textbf{the} \textbf{blue} is \textbf{the} prediction, how well it can predict, is to within \textbf{the} \textbf{width} of an atom on average, is how accurate \textbf{the} prediction is, which is what is needed for biologists to use it, and for drug design and for disease understanding, which is what AlphaFold \textbf{unlocks}.
CA : You made a surprising decision, which was to give away \textbf{the} actual results of your 200 million proteins.
DH : We open-sourced AlphaFold and gave everything away on a huge database with our wonderful colleagues, \textbf{the} European Bioinformatics Institute.
CA : I mean, you ' re part of Google.
Was there a phone call saying,"Uh, Demis, what did you just do?"DH : You know, I ' m lucky we have very supportive, Google ' s really supportive of science and understand \textbf{the} benefits this can bring to \textbf{the} world.
And, you know, \textbf{the} argument here was that we could only ever have even scratched \textbf{the} \textbf{surface} of \textbf{the} potential of what we could do with this.
This, you know, maybe like a millionth of what \textbf{the} scientific community is doing with it.
There ' s over a million and a half biologists around \textbf{the} world have used AlphaFold and its predictions.
We think that ' s almost every biologist in \textbf{the} world is making use of this now, every pharma company.
So we ' ll never know probably what \textbf{the} full impact of it all is.
CA : But you ' re continuing this work in a new company that ' s spinning out of Google called Isomorph.
DH : Isomorphic.
CA : Isomorphic.
Give us just a sense of \textbf{the} vision there.
What ' s \textbf{the} vision?
DH : AlphaFold is a sort of fundamental \textbf{biology} tool.
Like, what are these 3D structures, and then what might they do in nature?
And then if you, you know, \textbf{the} reason I thought about this and was so excited about this, is that this is \textbf{the} beginnings of understanding disease and also maybe helpful for designing drugs.
So if you know \textbf{the} shape of \textbf{the} protein, and then you can kind of figure out which part of \textbf{the} \textbf{surface} of \textbf{the} protein you ' re going to target with your drug compound.
And Isomorphic is extending this work we did in AlphaFold into \textbf{the} \textbf{chemistry} space, where we can design chemical compounds that will bind exactly to \textbf{the} right spot on \textbf{the} protein and also, importantly, to nothing else in \textbf{the} body.
So it doesn ' t have any side effects and it ' s not toxic and so on.
And we ' re building many other AI models, sort of sister models to AlphaFold to help predict, make predictions in \textbf{chemistry} space.
CA : So we can expect to see some pretty dramatic health medicine breakthroughs in \textbf{the} coming few years.
DH : I think we ' ll be able to get down drug discovery from years to maybe months.
CA : OK.
Demis, I ' d like to change direction a bit.
Our mutual friend, Liv Boeree, gave a talk last year at TEDAI that she called \textbf{the}"Moloch Trap."\textbf{The} Moloch Trap is a situation where organizations, companies in a competitive situation can be driven to do things that no individual running those companies would by themselves do.
I was really struck by this talk, and it ' s felt, as a sort of layperson observer, that \textbf{the} Moloch Trap has been shockingly in effect in \textbf{the} last couple of years.
So here you are with DeepMind, sort of pursuing these amazing medical breakthroughs and scientific breakthroughs, and then suddenly, kind of out of left field, OpenAI with Microsoft releases ChatGPT.
And \textbf{the} world goes crazy and suddenly goes,"Holy crap, AI is..."you know, everyone can use it.
And there ' s a sort of, it felt like \textbf{the} Moloch Trap in action.
I think Microsoft CEO Satya Nadella actually said,"Google is \textbf{the} 800-pound gorilla in \textbf{the} \textbf{search} space.
We wanted to make Google dance."How...?
And it did, Google did dance.
There was a dramatic response.
Your role was changed, you took over \textbf{the} whole Google AI effort.
Products were rushed out.
You know, Gemini, some part amazing, part embarrassing.
I ' m not going to ask you about Gemini because you ' ve addressed it elsewhere.
But it feels like this was \textbf{the} Moloch Trap happening, that you and others were pushed to do stuff that you wouldn ' t have done without this sort of catalyzing competitive thing.
Meta did something similar as well.
They rushed out an \textbf{open-source} version of AI, which is arguably a reckless act in itself.
This seems \textbf{terrifying} to me.
Is it \textbf{terrifying}?
DH : Look, it ' s a complicated topic, of course.
And, first of all, I mean, there are many things to say about it.
First of all, we were working on many large language models.
And in fact, obviously, Google research actually invented Transformers, as you know, which was \textbf{the} architecture that allowed all this to be possible, five, six years ago.
And so we had many large models internally. \textbf{The} thing was, I think what \textbf{the} ChatGPT moment did that changed was, and fair play to them to do that, was they demonstrated, I think somewhat surprisingly to themselves as well, that \textbf{the} public were ready to, you know, \textbf{the} general public were ready to embrace these systems and actually find value in these systems.
Impressive though they are, I guess, when we ' re working on these systems, mostly you ' re focusing on \textbf{the} flaws and \textbf{the} things they don ' t do and hallucinations and things you ' re all familiar with now.
We ' re thinking, you know, would anyone really find that useful given that it does this and that?
And we would want them to improve those things first, before putting them out.
But interestingly, it turned out that even with those flaws, many tens of millions of people still find them very useful.
And so that was an interesting update on maybe \textbf{the} convergence of products and \textbf{the} science that actually, all of these amazing things we ' ve been doing in \textbf{the} lab, so to speak, are actually ready for prime time for general use, beyond \textbf{the} rarefied world of science.
And I think that ' s pretty exciting in many ways.
CA : So at \textbf{the} moment, we ' ve got this exciting \textbf{array} of products which we ' re all \textbf{enjoying}.
And, you know, all this generative AI stuff is amazing.
But let ' s roll \textbf{the} clock forward a bit.
Microsoft and OpenAI are reported to be building or investing like 100 billion dollars into an absolute monster database supercomputer that can offer compute at orders of magnitude more than anything we have today.
It takes like five gigawatts of energy to drive this, it ' s estimated.
That ' s \textbf{the} energy of New York City to drive a data center.
So we ' re pumping all this energy into this giant, vast brain.
Google, I presume is going to match this type of investment, right?
DH : Well, I mean, we don ' t talk about our specific numbers, but you know, I think we ' re investing more than that over time.
So, and that ' s one of \textbf{the} reasons we teamed up with Google back in 2014, is kind of we knew that in order to get to AGI, we would need a lot of compute.
And that ' s what ' s transpired.
And Google, you know, had and still has \textbf{the} most \textbf{computers}.
CA : So Earth is building these giant \textbf{computers} that are going to basically, these giant brains, that are going to power so much of \textbf{the} future economy.
And it ' s all by companies that are in competition with each other.
How will we avoid \textbf{the} situation where someone is getting a lead, someone else has got 100 billion dollars invested in their thing.
Isn ' t someone going to go,"Wait a sec.
If we used reinforcement learning here to maybe have \textbf{the} AI tweak its own code and rewrite itself and make it so [ powerful ], we might be able to catch up in nine hours over \textbf{the} weekend with what they ' re doing.
Roll \textbf{the} dice, dammit, we have no choice.
Otherwise we ' re going to lose a fortune for our shareholders."How are we going to avoid that?
DH : Yeah, well, we must avoid that, of course, clearly.
And my view is that as we get closer to AGI, we need to collaborate more.
And \textbf{the} good news is that most of \textbf{the} scientists involved in these labs know each other very well.
And we talk to each other a lot at conferences and other things.
And this technology is still relatively nascent.
So probably it ' s OK what ' s happening at \textbf{the} moment.
But as we get closer to AGI, I think as a society, we need to start thinking about \textbf{the} types of architectures that get built.
So I ' m very optimistic, of course, that ' s why I spent my whole life working on AI and working towards AGI.
But I suspect there are many ways to build \textbf{the} architecture safely, robustly, reliably and in an understandable way.
And I think there are almost certainly going to be ways of building architectures that are unsafe or risky in some form.
So I see a sort of, a kind of bottleneck that we have to get humanity through, which is building safe architectures as \textbf{the} first types of AGI systems.
And then after that, we can have a sort of, a flourishing of many different types of systems that are perhaps sharded off those safe architectures that ideally have some mathematical guarantees or at least some practical guarantees around what they do.
CA : Do governments have an essential role here to define what a level playing field looks like and what is absolutely taboo?
DH : Yeah, I think it ' s not just about -- actually I think government and civil society and academia and all parts of society have a critical role to play here to shape, along with industry labs, what that should look like as we get closer to AGI and \textbf{the} cooperation needed and \textbf{the} collaboration needed, to prevent that kind of runaway race dynamic happening.
CA : OK, well, it sounds like you remain optimistic.
What ' s this image here?
DH : That ' s one of my favorite images, actually.
I call it, like, \textbf{the} tree of all knowledge.
So, you know, we ' ve been talking a lot about science, and a lot of science can be boiled down to if you imagine all \textbf{the} knowledge that exists in \textbf{the} world as a tree of knowledge, and then maybe what we know today as a civilization is some, you know, small subset of that.
And I see AI as this tool that allows us, as scientists, to explore, potentially, \textbf{the} entire tree one day.
And we have this idea of root node problems that, like AlphaFold, \textbf{the} protein-folding problem, where if you could crack them, it \textbf{unlocks} an entire new branch of discovery or new research.
And that ' s what we try and focus on at DeepMind and Google DeepMind to crack those.
And if we get this right, then I think we could be, you know, in this incredible new era of radical abundance, curing all diseases, spreading consciousness to \textbf{the} stars.
You know, maximum human flourishing.
CA : We ' re out of time, but what ' s \textbf{the} last example of like, in your dreams, this dream question that you think there is a shot that in your lifetime AI might take us to?
DH : I mean, once AGI is built, what I ' d like to use it for is to try and use it to understand \textbf{the} fundamental nature of reality.
So do experiments at \textbf{the} Planck scale.
You know, \textbf{the} smallest possible scale, theoretical scale, which is almost like \textbf{the} resolution of reality.
CA : You know, I was brought up religious.
And in \textbf{the} Bible, there ' s a story about \textbf{the} tree of knowledge that doesn ' t work out very well.
Is there any scenario where we discover knowledge that \textbf{the} \textbf{universe} says,"Humans, you may not know that."DH : Potentially.
I mean, there might be some unknowable things.
But I think scientific method is \textbf{the} greatest sort of invention humans have ever come up with.
You know, \textbf{the} enlightenment and scientific discovery.
That ' s what ' s built this incredible modern civilization around us and all \textbf{the} tools that we use.
So I think it ' s \textbf{the} best \textbf{technique} we have for understanding \textbf{the} enormity of \textbf{the} \textbf{universe} around us.
CA : Well, Demis, you ' ve already changed \textbf{the} world.
I think probably everyone here will be cheering you on in your efforts to ensure that we continue to accelerate in \textbf{the} right direction.
DH : Thank you.
CA : Demis Hassabis.

% matched lemmas: array, artificial, biology, blue, calculate, chemistry, computer, enjoy, fantastic, learning, left, open-source, search, surface, technique, terrifying, the, theme, universe, unlock, width
\end{textsample}
